\section{短幅摆线和长幅摆线}
\markboth{18. 短幅摆线和长幅摆线}{18. 短幅摆线和长幅摆线}

当动圆 $C$ 沿着定直线 $l$ 滚动时，固定在圆 $C$ 的平面内但不在圆周上的一点 $M$ 所画出的曲线叫做变幅摆线。变幅摆线可以分成两类：当 $M$ 点固定在圆 $C$ 内部时，曲线称为短幅摆线（图 18.1，上）；当 $M$ 点在圆 $C$ 外部时，曲线称为长幅摆线，又叫做余摆线（图 18.1，下）。

设圆 $C$ 的半径是 $r$，$M$ 点到圆心 $C$ 的距离是 $a$。取定直线 $l$ 为 $x$ 轴，并使动圆位于 $x$ 轴上方。又选取原点的位置，使得当动圆的圆心移动到 $y$ 轴上时，$M$ 点运动到最低位置。

\begin{center}
  \begin{minipage}{\linewidth}
    \centering
    \begin{tikzpicture}[scale=0.82, >=Stealth, font=\footnotesize]
      % 上图：短幅摆线 (a < r)
      \begin{scope}[yshift=3.15cm]
        \def\r{0.76}
        \def\a{0.42}
        \def\th{1.75}
        \pgfmathsetmacro{\cx}{\r*\th}
        \pgfmathsetmacro{\cy}{\r}
        \pgfmathsetmacro{\mx}{\r*\th - \a*sin(\th r)}
        \pgfmathsetmacro{\my}{\r - \a*cos(\th r)}
        \pgfmathsetmacro{\angM}{deg(1.5*pi - \th)}

        \draw[->] (-1.1, 0) -- ({2*pi*\r + 1.35}, 0) node[below left] {$x$};
        \draw[->] (0, -0.22) -- (0, {2*\r + 0.45}) node[left] {$y$};
        \node[below left] at (0, 0) {$O$};

        \draw[very thick, curveblue, smooth, domain=-1.3:{2*pi+1.3}, samples=100]
          plot ({\r*\x - \a*sin(\x r)}, {\r - \a*cos(\x r)});

        \coordinate (O1) at (0, 0);
        \coordinate (C1) at (\cx, \cy);
        \coordinate (P1) at (\cx, 0);
        \coordinate (M1) at (\mx, \my);

        \draw[thick] (C1) circle (\r);
        \draw[thin] (O1) -- (M1);
        \draw[thin] (C1) -- (M1);
        \draw[thin] (C1) -- (P1);
        \draw (C1) ++(270:0.25) arc (270:\angM:0.25);
        \node at ($(C1)+({0.5*(270+\angM)}:0.40)$) {\scriptsize$\theta$};

        \fill (C1) circle (1.2pt) node[right, inner sep=1.5pt] {$C$};
        \fill (M1) circle (1.2pt) node[above, inner sep=1.5pt] {$M$};
        \fill (P1) circle (1.2pt) node[below] {$P$};
        \node[right, inner sep=1pt] at ($(C1)!0.55!(P1)$) {\scriptsize$r$};
        \node[above right, inner sep=0.5pt] at ($(C1)!0.42!(M1)$) {\scriptsize$a$};
      \end{scope}

      % 下图：长幅摆线 (a > r)
      \begin{scope}[yshift=0cm]
        \def\r{0.76}
        \def\a{1.15}
        \def\th{2.35}
        \pgfmathsetmacro{\cx}{\r*\th}
        \pgfmathsetmacro{\cy}{\r}
        \pgfmathsetmacro{\mx}{\r*\th - \a*sin(\th r)}
        \pgfmathsetmacro{\my}{\r - \a*cos(\th r)}
        \pgfmathsetmacro{\angM}{deg(1.5*pi - \th)}

        \draw[->] (-1.1, 0) -- ({2*pi*\r + 1.35}, 0) node[below left] {$x$};
        \draw[->] (0, -0.58) -- (0, {2*\r + 0.72}) node[left] {$y$};
        \node[below] at (-0.38, -0.02) {$O$};

        \draw[very thick, curveblue, smooth, domain=-2.15:{2*pi+2.15}, samples=130]
          plot ({\r*\x - \a*sin(\x r)}, {\r - \a*cos(\x r)});

        \coordinate (O2) at (0, 0);
        \coordinate (C2) at (\cx, \cy);
        \coordinate (P2) at (\cx, 0);
        \coordinate (M2) at (\mx, \my);

        \draw[thick] (C2) circle (\r);
        \draw[thin] (O2) -- (M2);
        \draw[thin] (C2) -- (M2);
        \draw[thin] (C2) -- (P2);
        \draw (C2) ++(270:0.28) arc (270:\angM:0.28);
        \node at ($(C2)+({0.5*(270+\angM)}:0.44)$) {\scriptsize$\theta$};

        \fill (C2) circle (1.2pt) node[right, inner sep=1.5pt] {$C$};
        \fill (M2) circle (1.2pt) node[above left, inner sep=1.5pt] {$M$};
        \fill (P2) circle (1.2pt) node[below] {$P$};
        \node[right, inner sep=1pt] at ($(C2)!0.55!(P2)$) {\scriptsize$r$};
        \node[above right, inner sep=0.5pt] at ($(C2)!0.42!(M2)$) {\scriptsize$a$};
      \end{scope}
    \end{tikzpicture}
    \captionof*{figure}{图 18.1}
  \end{minipage}
\end{center}

利用图 18.1 容易求得变幅摆线的参数方程是
\[
  x = rt - a\sin t,\quad y = r - a\cos t.
\]
当 $a < r$ 时得到短幅摆线，$a > r$ 时得到长幅摆线。又若令 $a = r$，就回到普通摆线的方程。

由于 $|\cos t| \leqslant 1$，所以当 $a < r$ 时恒有 $y > 0$，即短幅摆线与 $x$ 轴无公共点。当 $a > r$ 时 $y$ 的值可以为零，也可以为负，这意味着长幅摆线与 $x$ 轴相交，并且有一部分弧出现在 $x$ 轴下方\footnote{重排校订：原书此处误印作“$y$ 轴下方”，据上下文及参数方程 $y = r - a\cos t$ 订正为“$x$ 轴下方”。}。% 原书此处误印作“y 轴”
从图 18.1 看出，长幅摆线自己绕成许多小圈，叫做绕扣，每个绕扣都被 $x$ 轴分成上下两截，并且长幅摆线与 $x$ 轴的任何交点都在绕扣上。

在日常生活中很容易找到变幅摆线的实例。还是以自行车为例，当一个人推着自行车慢慢往前走时，旁观者就可以看到车轮钢圈内侧的彩色气门帽慢慢地描绘出一条短幅摆线。
